\chapter{Literature Survey} \label{chapter2}

\section{Scope of the Survey}

This chapter reviews the literature that is load-bearing for the V3 design: (i) the directional forecasting problem and the random-walk null; (ii) supervised learning for tabular financial features (gradient boosting); (iii) deep sequence models for daily OHLCV; (iv) ensembling and meta-learning; (v) probability calibration; (vi) backtest hygiene and lookahead bias; (vii) macro / exogenous cues relevant for the NSE; (viii) Indian-market specific studies. Each subsection ends with a short note on the design implications for V3.

\section{The Directional Forecasting Problem and the EMH Null}

The semi-strong form of the Efficient Market Hypothesis~\cite{fama1970emh} implies that, conditional on public information, expected next-day excess returns are zero. Empirically, however, the literature shows that the joint hypothesis of zero predictability is rejected for short horizons under specific conditioning sets: weekly momentum~\cite{jegadeesh1993momentum}, post-earnings drift, market-microstructure inefficiencies, and macro-cue spillover from the US session into the Asian open. The V3 task is therefore not whether prices are predictable in general but whether \emph{a measurable subset of trading days carries enough conditional structure} for a learnt classifier to beat the round-trip cost.

\textbf{Implication for V3.} The gate is set so that the model is only allowed to trade on days where calibrated probability exceeds $0.58$ \emph{and} the predicted absolute move exceeds the round-trip cost ($\sim 0.4\%$ for NSE delivery). The reported headline number on all days (51.25\%) is consistent with weak-form efficiency; the headline number on signal days (60.33\%) is the figure that matters operationally.

\section{Gradient Boosting for Tabular Finance}

\textbf{XGBoost}~\cite{chen2016xgboost} introduced regularised tree boosting with second-order Taylor optimisation, sparsity-aware split-finding, and weighted quantile sketches. \textbf{LightGBM}~\cite{ke2017lightgbm} replaced level-wise growth with leaf-wise growth, gradient-based one-side sampling, and histogram-based feature bucketing, lowering wall-clock cost on wide tabular data without surrendering accuracy. On panel-style daily OHLCV with hand-engineered features, gradient boosting variants consistently outperform un-engineered deep nets~\cite{shwartz2022tabular,grinsztajn2022treesvsdl}. Both libraries handle missing values natively and tolerate heterogeneous feature scales without normalisation, which is convenient when ratios, percentages and raw indices coexist in the same matrix.

\textbf{Implication for V3.} V3 trains LightGBM and XGBoost on the post-PCA scaled matrix, with explicit $\ell_1$ and $\ell_2$ regularisation ($\alpha=0.3$, $\lambda=1.5$), a 1{,}000-tree budget, learning rate $0.01$, depth 5, and early stopping after 50 non-improving rounds. These two trees provide the meta-learner's feature space.

\section{Deep Sequence Models on Daily OHLCV}

The LSTM~\cite{hochreiter1997lstm} and the GRU~\cite{cho2014gru} were the first sequence architectures with controlled long-range gradient flow. Their direct application to financial time series~\cite{fischer2018lstm,sezer2020survey} produces respectable equity-curve numbers but mediocre out-of-sample direction accuracies once leakage is removed. Bidirectional variants are admissible here because every feature is lagged (no peek through Bi-direction). Convolutional preprocessing (CNN--LSTM, CNN--GRU) lifts local pattern extraction off the recurrent stack and is causal-padded to preclude future-step leakage. Temporal Convolutional Networks~\cite{bai2018tcn} replace the recurrence with dilated causal convolutions and have receptive field $1 + (k-1) \cdot 2 \sum d_i$, which for kernel size 3 and dilations $\{1, 2, 4, 8\}$ is $\sim 61$ time steps, comfortably wider than the 20-step sequence we feed in. N-BEATS~\cite{oreshkin2020nbeats} introduced residual backcast blocks for univariate forecasting; we adapt the generic configuration with a sigmoid classification head. Lim et al.'s Temporal Fusion Transformer~\cite{lim2021tft} influenced our \texttt{TCN-Transformer} block, which composes TCN local feature extraction with multi-head self-attention.

\textbf{Implication for V3.} Eight deep-learning models are trained per symbol per window with identical input shapes $(B \times 20 \times d)$ (batch size $B$, sequence length 20, feature dimension $d$), fed from the robust-scaled matrix. They consume the same input as the trees but contribute independently to the meta-learner, providing decorrelated error structure rather than competing for raw accuracy.

\section{Ensembling, Stacking, and the Meta-Learner}

\textbf{Stacked generalisation}~\cite{wolpert1992stacking} feeds out-of-fold predictions of base learners into a level-2 model that learns the optimal blend. The level-2 model should be simple to avoid further overfit; logistic regression is the canonical choice for binary stacking. The variance-reduction case for heterogeneous ensembles is made by Dietterich~\cite{dietterich2000ensemble}, who distinguishes bagging (same model, different data), boosting (sequential error-correction) and stacking (different model classes, learned blend). Heterogeneity is the only one of the three that gives V3 a structurally different error signal across families; bagging is implicit inside each booster, boosting is implicit inside each booster, and stacking is the V3 layer.

\textbf{Implication for V3.} V3 stacks the two tree probabilities (and optionally the deep-learning probabilities, when their validation accuracy clears a sanity floor) via logistic regression on out-of-fold tree predictions on the validation slice, producing a combined meta-probability at inference.

\section{Probability Calibration}

A classifier can rank the right way and still mis-state its probabilities. Modern deep networks in particular are systematically over-confident~\cite{guo2017calibration}. Two corrections dominate the literature: Platt scaling (a single logistic regression on the validation set's logits) and isotonic regression (a non-parametric monotonic fit). \textbf{Temperature scaling}~\cite{guo2017calibration} is the one-parameter limit of Platt scaling: $\hat p = \sigma(z / T)$ with $T > 0$ fit by NLL on the validation set. It preserves the argmax (and therefore the accuracy) while contracting/expanding the probability distribution to match observed frequencies.

\textbf{Implication for V3.} A temperature scalar $T^\star$ is fit per fold per symbol on the validation slice by one-dimensional scalar optimisation over $T \in [0.05,\ 5.0]$ minimising binary negative log-likelihood. It is then applied to the meta-probability at inference. Without this step, the $0.58$ gate threshold has no operational meaning; with calibration, $0.58$ corresponds to an empirically near-reliable probability.

\section{Backtest Hygiene and Lookahead Bias}

Bailey, Borwein, L\'opez de Prado, and Zhu~\cite{bailey2014lookahead} and L\'opez de Prado~\cite{prado2018walkforward} cataloguing the failure modes of financial backtests: random train/test splits leak future information; full-panel scalers leak the test mean into the train; rolling statistics computed before the split leak forward; full-panel feature selection leaks target signal; and reporting only the best of many configurations under-states the multiple-testing penalty. Pesaran and Timmermann~\cite{pesaran1995predictability} establish that the apparent predictability of stock returns is fragile to the conditioning set.

\textbf{Implication for V3.} Every audit point in~\cite{prado2018walkforward} maps to a concrete control in the pipeline (cf.\ Chapter~\ref{chapter3}, the leakage audit table): the target label is built by a single forward shift applied \emph{after} all features are computed; scalers are refit per fold on the training slice only; the global cues are merged using a strict prior-day backward temporal join; feature selection is restricted to the training window; only one canonical run is reported and its configuration is committed with fixed random seeds.

\section{Macro Cues and the Indian Calendar}

Two strands of literature matter for the Indian session. First, \textbf{global cues}: the US-close-to-Asian-open spillover is well documented~\cite{eun1989international,rapach2013}; in V3 we add S\&P 500, Nasdaq, US VIX, DXY, Crude and Nikkei, all shifted by one calendar day, on top of Nifty 50 and Nifty Bank. Second, the \textbf{NSE microstructure calendar}: F\&O monthly expiry (last Thursday of each month, with weekly Bank-Nifty options on Wednesdays / Thursdays depending on the cohort), RBI MPC announcement days, the Union Budget date, and the quarterly results season~\cite{nse_handbook}. Each of these compresses or expands the conditional variance of the next-day move and creates short, predictable windows that a learner can pick up if given the calendar.

\textbf{Implication for V3.} V3 force-includes a global market-cue feature set (S\&P 500, Nasdaq, VIX, DXY, Crude, Nikkei returns) for every stock, additionally includes USD/INR and Nasdaq features for IT-sector stocks, and banking-sector cue features (including days to the next RBI MPC meeting) for banking stocks. Calendar features (days to the nearest F\&O expiry, RBI meeting, Union Budget, and a results-season indicator) are computed from fixed date lists.

\section{Indian-Market Prior Work}

A representative sample of Indian-market forecasting papers, with the headline accuracy and the protocol that produced it, appears in Table~\ref{tab:indian_priors}. The pattern is consistent: small universes (5--20 stocks), small feature sets (10--50 features), random splits or single chronological splits, and headline accuracies in the 56--65\% range. None reports a panel-wide expanding-window walk-forward with regime-conditional Sharpe.

\begin{table}[H]
\centering
\caption{Representative prior work on supervised equity-direction forecasting}
\label{tab:indian_priors}
\renewcommand{\arraystretch}{1.15}
\footnotesize
\begin{tabular}{@{}l l l c l@{}}
\toprule
\textbf{Study} & \textbf{Models} & \textbf{Universe} & \textbf{Feats.} & \textbf{Protocol / Acc.} \\
\midrule
Shah et al.~\cite{shah2019taxonomy}    & Review (taxonomy)       & survey          & ---  & Systematic review; 56--83\% reported across studies \\
Patel et al.~\cite{patel2015nse}       & ANN, SVM, RF, NB        & NSE, BSE        & 10   & Single split; up to 83\% (direction) \\
Long et al.~\cite{long2019deep}        & CNN-LSTM                & A-shares (CSE)  & 14   & Single split; $\sim$58\% (direction) \\
Nabipour et al.~\cite{nabipour2020deep}& XGBoost, LSTM           & Tehran exchange & 10   & Single split; variable by horizon \\
\midrule
V3 (this thesis)                       & 5-model + temp.\ cal.   & 100 Nifty-100   & $\sim$200 & 6-fold WFV; 51.25\% (all) / \textbf{60.33\%} (gated) \\
\bottomrule
\end{tabular}\\[4pt]
\footnotesize\itshape Gated row: calibrated probability $\ge 0.58$ and predicted move $\ge 0.4\%$; bootstrap CI $[57.99\%,\ 62.81\%]$.
\end{table}

\section{What V3 Inherits and What It Changes}

V3 inherits from the literature: gradient boosting on hand-engineered tabular features (Section~3); causal CNN preprocessing for recurrent stacks; the stacked-generalisation principle for the meta-layer; temperature scaling for probability outputs; expanding-window walk-forward as the validation protocol; and the global-cue / calendar conditioning set. V3 changes: it does not report a single chronological split as the headline; it does not stack untemperature-scaled probabilities; it does not use full-panel scalers; it does not skip the brokerage model in the backtest; and it does not promote a model to live capital on the basis of out-of-sample paper accuracy alone; the promotion gate also requires Brier-drift, slippage, and fill-rate evidence to have accumulated.
